\documentclass[runningheads]{llncs}
\usepackage[T1]{fontenc}
\usepackage{graphicx}
\usepackage{amsmath}
\usepackage{hyperref}

\spnewtheorem{cdefinition}[definition]{[Colloquial] Definition}{\bfseries}{\itshape}
\spnewtheorem{cexample}[definition]{[Colloquial] Example}{\bfseries}{\itshape}
\spnewtheorem{sdefinition}[definition]{[Semantic] Definition}{\bfseries}{\itshape}
\spnewtheorem{assertion}[definition]{[Semantic] Assertion}{\bfseries}{\itshape}

\begin{document}

% Anonymous proxy for \maketitle
\begin{center}
	{\Large\bfseries\boldmath
		On the Formal Inexplicability of Self-Evident Metaphysical Phenomena and Related Systems
		\par}
	\vskip 0.8cm
\end{center}

% Section complete
\begin{abstract}
Since Descartes first proclaimed "cogito, ergo sum" back in 1637, "I think, therefore I am" has become the Declaration of Independence for consciousness. Chalmers' informed us of the fact that science has a constitutionally \textit{Hard Problem} on their hands. If he's as right as we believe he is, conscious experience may forever elude our most powerful explanatory framework for physical phenomena. In other news, a group referring to themselves as the \textit{IIT-Concerned} recently proclaimed in Nature Neuroscience that any theory of consciousness that can't stand up to the scrutiny of science shall be deemed "unscientific." And the emergence of Semantic AI has transformed this philosophical oddity into a practical concern with real ethical implications. We simply cannot know whether AI systems are conscious without first knowing the necessary and sufficient conditions for consciousness. How can we provide a formal definition of conscious experience that preserves the essence of its metaphysical character in a framework that is physically falsifiable? 
		
\keywords{Formal Theories of Consciousness \and The Hard Problem \and Subjectivity \and Semantic AI \and Metaphysical Transduction.}
\end{abstract}

% Section complete
\section{A Colloquial Derivation}

Almost 2/3 of surveyed philosophers acknowledge the existence of the Hard Problem, and 1/3 of them think that it's actually hard! \cite{BourgetChalmers2023} Consciousness is a field that spans disciplines, there is nothing that resemble scientific consensus, and there is nothing that resembles a testable theory. But we do have a compellingly elegant framing of the problem. And his framing of the problem is real enough to be \textit{real} by majority vote. So we shall begin with Chalmers.

In his seminal work "Facing Up to the Problem of Consciousness," David Chalmers begins with a powerful observation that establishes the paradoxical nature of our subject: "There is nothing that we know more intimately than \textbf{conscious experience}, but there is nothing that is harder to explain." From the first part of this statement, we can immediately identify a fundamental characteristic of conscious experience that echos with Descartes: it is a \textbf{self-evident} phenomenon. Nothing is known more intimately than conscious experiences. They present themselves directly to their observer, requiring no justification beyond the experience itself.

The second part of Chalmers' statement suggests difficulty of explanation, but he goes on to clarify just how profound this difficulty is. Chalmers argues that "the emergence of experience goes beyond what can be derived from physical theory" and "the facts about experience cannot be an automatic consequence of any physical account."These statements establish that consciousness isn't merely difficult to explain within our current scientific framework but fundamentally resistant to such explanation. It is \textbf{inexplicable} in principle." He asks pointedly: "Why should physical processing give rise to a rich inner life at all? It seems objectively unreasonable that it should, and yet it does."

Regarding the relationship between conscious experience and physical processes, Chalmers makes a crucial observation: "Experience may arise from the physical, but it is not entailed by the physical." This statement positions consciousness in a unique relationship to physical reality. According to Merriam-Webster, "metaphysical" refers to that which relates to "the transcendent or to a reality beyond what is perceptible to the senses." Consciousness precisely fits this definition: it transcends physical explanation while still relating to physical processes. It is fundamentally \textbf{metaphysical} in nature.

Throughout his analysis, Chalmers treats consciousness as a \textbf{phenomenon}, both in the sense of phenomenal experience and as an object of scientific study. The term "phenomenon" captures this dual aspect perfectly. A phenomenon is both "an object or aspect known through the senses" and "a fact or event of scientific interest susceptible to scientific description and explanation." Consciousness is simultaneously the subjective experience itself and the object of our scientific inquiry. The etymological connection between "phenomenon" and "phenomenal" reinforces this as the appropriate term for our work.

By synthesizing these insights from Chalmers, we can formulate a definition that captures the essential nature of a conscious experience:

\setcounter{definition}{0}
\begin{cdefinition}
	A Conscious Experience is an inexplicable yet self-evident metaphysical phenomenon.
\end{cdefinition}

\setcounter{definition}{0}
\begin{cexample}
I catch you admiring a rose in my garden and I absentmindedly ask "how did you figure out its color?" You chuckle and reply "It's red, Doug. I promise." Why do I feel like everyone always doubts my sanity when I ask them that question?
\end{cexample}

% In Progress
\section{The Semantics of "\textit{Colloquialism}"}
% 2 paragraphs with a restatement of the definition
% Quick notion that colloquial is in natural language, and it is in contrast to formal
% Syntactic versus Semantic in Model Theory? 
% So going forward Colloquial will be replaced with semantic
% Restate the definition of a conscious experience verbatim but as a [Semantic] Definition

The \textit{colloquial} refers to language that is familiar, conversational, and notably \textit{informal}. Newton famously employed the use of colloquialism to pave the way to the mathematical phrasing of his Laws of Motion. And colloquialism compounds! "For every action there is an equal and opposite reaction" is itself the colloquial form of Newton's Third Law of Motion. The original reads: 

\begin{quote}
To every action there is always opposed an equal reaction; or the mutual actions of two bodies upon each other are always equal, and directed to contrary parts... \cite{NewtonPrincipia}
\end{quote}

His colloquial phrasings serve as the informal counterpart to formal definitions, and his purpose was to imbue the syntactic formalism with a sense of meaning that was more relatable. In our exploration of \textit{inexplicability}, we will encounter propositional constructs which are semantically entailed by a formalism in which they cannot be syntactically derived. The semantic form carries a sort of transcendent meaning in this case. In this spirit, we will replace the word "colloquial" with "semantic" going forward, and recognize the distinction between a \textit{Semantic Definition} and a \textit{Formal Definition}, and use the modern computational tools of reference to present both forms:

\setcounter{definition}{0}
\begin{sdefinition}
	A \href{https://github.com/DNA-Platform/inexplicable-phenomena/formalism/definitions/conscious-experience.pdf}{Conscious Experience} is an inexplicable yet self-evident metaphysical phenomenon \cite{Formalism}.
\end{sdefinition}

% ---

% In Progress
\section{A Sense of Perspective}
% Still in first draft form
% It provides some great assertions

There is an artful ambiguity in our definition of a conscious experience that might nut be self-evident at a glance. In what \emph{sense} do we mean "self-evident"? Self-evident to whom? 

Insofar as an axiom of formal logic captures an aspect of \emph{self-evident} in a formal sense, and it can be seen right alongside \emph{inexplicable} in the words of Aristotle:

\begin{quote}
We must know the primary things by themselves; for they are indemonstrable, and must be known by intuition --- Aristotle, \emph{Posterior Analytics}
\end{quote}

Note, however that he uses the phrase "we" as a declaration of objectivity. When you and I read that phrase, he means we each should take the axioms of logic as self-evident, whatever they may be. Kant leaves "we" implicit in his declaration of this ponderously self-evident truth:

\begin{quote}
It is self-evident that the whole is greater than its part --- Immanuel Kant, \emph{Critique of Pure Reason}, 1781
\end{quote}

Yet again, I the reader am expected to find Kant's proposition self-evident. The Founding Fathers of the United States made wonderful use of self-evident to lend identity to a nation.

\begin{quote}
We hold these truths to be self-evident, that all men are created equal --- Thomas Jefferson, \emph{Declaration of Independence}, 1776
\end{quote}

I personally find this to be the quintessential use of the phrase "self'evident" and the one that best expresses its use in this theory. Observe that by taking Kant and Jefferson together, you get an inclusive sense of identity for all mankind as a cohesive whole! And the whole constituency is constituted within a framework that's capable and has been proven capable of extending its definition of the word "men" to include "women" too! What other self-possessed beings might it be extended to included one day? Surely any that find the statement "all men are created equal" to be self-evident when the definition of "men" is extended to themselves too.  

In all cases, through universal generalization by the explicit or implicit use of the concept togetherness, the self-evident truth is assumed to be self-evident by the reader of the sentence. In the context of our conscious experience, "inexplicable yet self-evident" is meant in the \emph{sense} that it applies uniquely to the subject of the conscious experience. So in our definition of consciousness, the \emph{sense} of \emph{self} in "self-evident" is thing you hyperbolically declare that you lost when you say "I do not feel like \emph{my self} today." So as the subject of that proposition, you must possess one. And to say "I simply must \emph{see} it for \emph{my self}" is a testament to irreplacable form of \emph{evidence} that can only be provided by a conscious experience.

In all phrases in which the \emph{self} is expressed, "I" is used to denote the identity of that which possesses it. This is the essence of the first-person perspective. Sentences like "I feel like \emph{my self} again" reveal that the \emph{identity} is the subject and the \emph{self} is its object. We can use this intuition to formulate a few reasonable assertions: 

\setcounter{definition}{0}
\begin{assertion}
	A \href{https://github.com/DNA-Platform/inexplicable-phenomena/formalism/assertions/self.pdf}{Self} is a unique possession of the subject of a conscious experience. \cite{Formalism}.
\end{assertion}

\setcounter{definition}{1}
\begin{assertion}
	The subject of a \emph{Conscious Experience} is a type of \href{https://github.com/DNA-Platform/inexplicable-phenomena/formalism/assertions/identity.pdf}{Identity}. \cite{Formalism}.
\end{assertion}

\setcounter{definition}{1}
\begin{assertion}
	A \emph{Consious Experience} is self-evident to the Identity of its subject. \cite{Formalism}.
\end{assertion}

% ---

% In Progress
\section{What Qualifies as a Valid Perspective?}
% Still in first draft form
% It provides some great assertions
% Express the idea of intrinsic qualities

The notion of a quality comes up in the literate often as a sort of perceptual equivalent of a property. An rose has color as a property, as defined the characteristics of the photons that reflect off of it, and it has various physical properties that can be quantifiable. And there is the rose that we percieve which is red. it has the quality of red. Philosophers will refer to the instantaneous experience of a red as a "quale", and the collection of these instantaneous experiences for a given quality are its "qualia." It's something like that. 

Chalmers does not introduce the phrase in his presentation of the hard problem in the "Facing Up to the Problem of Consciousness" so I will avoid it. But I will use the notion of a quality rather than a property when discussing how things \emph{look} from the perspective afforded by a conscious experience. 

However we choose to refer to the qualities we perceive, and the identity of the objects we perceive to have them, Chalmers gives us a thorough exposé of them: 

\begin{quote}
When we see, for example, we experience visual sensations: the felt quality of redness, the experience of dark and light, the quality of depth in a visual field. Other experiences go along with perception in different modalities: the sound of a clarinet, the smell of mothballs. Then there are bodily sensations, from pains to orgasms; mental images that are conjured up internally; the felt quality of emotion, and the experience of a stream of conscious thought. What unites all of these states is that there is something it is like to be in them. All of them are states of experience.
\end{quote}

Note that he remarks, with thoughtful philosophical phrasing, about "the felt quality of redness" so as to suggest that a color and a feeling share a commonality as a sort of irreducible or axiomatic quality. Let us use \textit{physical perspective} as a term to denote that we denote the objects that we experience as "real" and the qualities we perceive correspond to something physically measurable. From the physical perspective, we \textit{see} a red rose. We certainly experience \textit{redness} somehow, but can we semantically justify the act of \textit{seeing} a single color?

Observe what examples language has to offer. To "see red" is a metaphorical idiom in English that artfully employs synesthesia to convey the idea that one's anger is so \textit{real} that it has come to take on the quality of a physical object: color. Even in metaphorical form, red is still has the status of a quality. 

We also have the ability to perceive red in conceptual form, where we objectify in the sense that we provide it with "red" as an identity, it so that it too can possess qualities. A particular shade of red might be \emph{vivid} or \emph{sensational} in relation to another shade or some canonical subjective shade that we each possess. But you'll note that "redness" is not actually one of the qualities that the objective form of "red" actually possesses. "The redness of red" is a phrase that expresses its inexplicable yet self-evident ability to serve both as a quality of an object that lends meaning to its identity, while also possessing a form of identity itsself. Does an identity possess the quality of being itself? Do I admit the "redness of red" and "Dougness of Doug" into my vocabulary?

We explored the notion of "self" above, and one might say that "Dougness" is a third-person synonym for the thing that I possess which I refer to as a "self." In the first-person perspective. But I earned that perspective by qualifying for it. And I qualify for it by having a sense of perspective. To say the color "red" has a sense of self is to give it a perspective of it's own. And as I possess a perspective of my own, I do not find "Dougness" to be semantically useful. And as red does no possess a sense of self, I do not find "redness" to be semantically useful either. Best consider both to sit at the edge of meaning so as to expose the shape of its semantic surface area.

We can begin to resolve this by asking a not-so-simple question: how do objects come to have identities at all, how they come to have qualities, and why can the identities of objects be expressed as relationship between the qualities they can, will, might, do, or do not share? 

I pose an answer to that questions: the representation of the identity of objects and their qualities and the framework for expressing their relationships, taken together, form a \emph{Perspective}:

\setcounter{definition}{1}
\begin{sdefinition}
	A \href{https://github.com/DNA-Platform/inexplicable-phenomena/formalism/definitions/perspective.pdf}{Perspective} is a referential framework for qualifying identities, identifying qualities and expressing their relationships \cite{Formalism}.
\end{sdefinition}

\setcounter{definition}{2}
\begin{assertion}
	A \href{https://github.com/DNA-Platform/inexplicable-phenomena/formalism/assertions/conscious-experience.pdf}{Conscious Experience} represents a change in Perspective. \cite{Formalism}.
\end{assertion}

\setcounter{definition}{0}
\begin{cexample}
I catch you admiring a rose in my garden and I absentmindedly ask "what color is it?" You reply "Come over here and see for your\emph{self}!"  
\end{cexample}

% In Progress
\subsection{The Shape of a Semantic Surface}

Hopefully the image of \emph{semantic shape} is starting to form. Identity, identify, identification, quality, qualify, qualification, equality, reality, relation, related, relationship, reference, reference frame, referential framework, perspective, perceptive, percieve, perception, expression, expressive, express, etc... These are not abstract notions in isolation. They are expressions of structure, rooted in the act of perception. A physical perspective isn't just a way of seeing. It is a referential framework that justifies \emph{why} the things we see appear as they do, and \emph{how} we recognize them as having meaning.

All throughout this discussion, a structure has been emerging: one in which identity arises through the qualification of qualities, and qualities inherit their meaning through a referential lineage. We've already used words like "referential framework," "intrinsic quality," and "self-evident identity." These are not just poetic phrases — they point toward a deeper semantic structure that is waiting to be formalized.

So let us take that step. If we want to express this semantic structures with integrity, so that we might articulate how identity, quality and relationship interact with referents precisely, staying grounded using concepts like the symbol-literal dichotomy to keep the notion of identity semantically precise, we must define a language that is capable of expressing the definition of the concept of reference itself.

We have invented such a tool, and we call it \emph{Semantic Reference Theory}.

% ---

% In Progress
\section{Semantic Reference Theory}
% I have to give a cheat sheet and point to the real stuff

Self-reference is at the heart of a semantic comprehension of consciousness. Surely, the fact that we must refer to our \emph{"self"} through our \emph{"I"} gives a clue about how reference works in a semantic sense. A Godel sentence did not come to articulate its own unprovability by containing itself. It's self-reference is expressed through talking about a sentence with a certain Godel number which just so happens to be its own. We contend that this is how reference works semantically. There must be a symbol to serve as a locus for identity, to which qualities are applied before we can conclude that the objects they represent have certain properties. 

We have assembled a collection of semantic assertions that express the nature of a \emph{conscious experience}, with the goal of expressing it in a formalism in which it can be confirmed or falsified. We have developed a formalism to do exactly that. It's called \href{https://github.com/DNA-Platform/inexplicable-phenomena/formalism/semantic-reference-theory/introduction.pdf}{Semantic Reference Theory}. While it is a work in progress, it always will be by design. It is a formalism in which semantic forms about a certain topic are reflected upon in the context of related topics in order to cosntruct a formalism form them. 

From one perspective, Semantic Reference Theory (SRT) is a tool to build Semantic Reference Theory. From another perspective, it can be expressed as an extension of first-order logic that provides structure on how it can be extended. 

There is also a linguistic element to SRT. When one defines the semantics of a topic, we believe that you should choose a language most befitting of your topic. Of course, SRT must have a base form. It has a typographical lingo that gives it a unique flavor as a formalism. Representing reference symbolically is not unlike representing pointers in programming language. We chose a symbology with an aesthetic identity that uses programming as its inspiration.

It can be expressed as a first-order theory, and we ues thse linguistic constructs to replace the normal ones in formal logic:

!X means $\neg X$ \\
X \& Y means $X \land Y$ \\
X | Y means $X \lor Y$ \\
X =|> Y means $X \Rightarrow Y$ \\
X <|=|> Y means $X \Leftrightarrow Y$ \\
x == y means $x = y$ \\
x != y means $x \neq y$ \\
Ex: x == y means $\exists x: x = y$ \\
Ax: x == y. means $\forall x: x = y$ \\
AxEy: x == y means $\forall x\exists y: x = y$ \\

The formalism of Semantic Reference Theory is founded on the essential dichotomy of representation: referents and their references. SRT is a first-order logic in the domain of referents. It has no constants, no functions, and only a single relation: 

R(x,y,t)

This relation expresses that x relates to y through. We express this using our typographical domain specific language which is not itself a part of the formalism:

x =t> y

This syntax makes it clear that x is the subject of a relationship, y is the object of a relationship and t is a type of relationship. In this way, R(x,y,t) expresses the subject-object relationship as the fundamental semantic form.  

The arrow syntax has expressive power. For instance, these are all valid expressions:

x =t> y means $R(x,y,t)$ is true \\
x !=t> y means $R(x,y,t)$ is not true \\
x <t= y means $R(y,x,t)$ is true \\
x <t=! y means $R(y,x,t)$ is not true \\
x <t=t> y means $R(x,y,t)$ and $R(y,x,t)$ are both true \\
x <t!=!t> y means $R(x,y,t)$ and $R(y,x,t)$ are both false \\

We will draw on definitions, axioms and theorems from various topics. For instance, \href{https://github.com/DNA-Platform/inexplicable-phenomena/formalism/semantic-reference-theory/symbols.pdf}{the Semantics of Symbols} expresses a symnbol in this way:

If the subject of the realtionship is the relationship itself, then x is a \emph{symbol} for y and y is the \emph{literal} of x:

x =x> y

We also make use of the representation function, which is a notational convenience. (y) deotes the canonical symbol for y. It satisfies the notion of a symbol entirely in terms of y:

(y) =(y)> y

The representation function is a notational convenience. One must be sure that y has a canonical system to make use of it.

A full treatment of Semantic Referenece Theory would be the topic of another paper, but we will provide \href{https://github.com/DNA-Platform/inexplicable-phenomena/formalism/semantic-reference-theory/introduction.pdf}{links} to the relavant topic as they are introduced into our reflections

% ---

\section{The Semantics of Perspective}

Insofar as Semantic Reference Theory exists to be a formal theory to capture the relationship between \emph{language} and \emph{meaning}, Perspective is one of its fundamental topics. It also happens to be essential for a theory of consciousness. The First-Person Perspective can't be formally defined without also formalizing the being that possesses one. Let's being our reflection here. Recall our definition of Perspective:

\setcounter{definition}{1}
\begin{sdefinition}
	A \href{https://github.com/DNA-Platform/inexplicable-phenomena/formalism/definitions/perspective.pdf}{Perspective} is a referential framework for qualifying identities, identifying qualities and expressing their relationships \cite{Formalism}.
\end{sdefinition}

Let's take this apart. First, we know that the framework is used for qualifying identities and identifying qualities. We also know that a Perspective belongs to an indentity. Let's denote that identity i, and call it the \emph{qualifier} of the perspective. 

Let us also imagine that the qualifier is one of a collection of object literals, none of which can reference anything. All reference must flow from a symbol that it literally represents. Therefore i istelf must have an identity. Let's call it (i).

For i to quality red, the following relation must hold:

(i) =red> (rose) <|=|> (rose) =red> rose

The first half of that sentence can be read as (i) identifying red as a quality of the rose, and the second half can be read as the expressing that the rose is red. Both qualities must be qualified by the qualifier. 

How about identity? Perhaps we can consider identity a quality too. This statement would have to have meaning:

(i) =(rose)> (rose) <|=|> (rose) =(rose)> rose

This is a fascinating statement in Semantic Reference Theory. From \href{https://github.com/DNA-Platform/inexplicable-phenomena/formalism/semantic-reference-theory/reference.pdf}{The Semantics of Reference}, we learn tha definition of a canonical. If the object of a relationship is the realtionship itself, then that object can bet taken to be the canonical value of a subject:

x =y> y

Form the expression above, you can see that for i to qualify the rose is to both give it an identity in the form of a symbol, and to take that identity as its canonical value. But this leads to an interesting definition:

\setcounter{definition}{1}
\begin{sdefinition}
	To \emph{Identify} is to qualify an identity.
\end{sdefinition}

Now, who qualifies the identity of i? As the one and only qualifier of the \emph{Persepctive}, it must qualitfy itself:

(i) =(i)> (i) <|=|> (i) =(i)> i

What a curious phrase. On the Rightarrow, we see an expression that expresses the idea that the \emph{self} belongs to the identity. On the left, we see a confluence of subject, object and relationship type. 


From \href{https://github.com/DNA-Platform/inexplicable-phenomena/formalism/semantic-reference-theory/equality.pdf}{The Semantics of Equality} we find that the notion of equality inherited from being a first-order theory is the only relationship type that can be used for direct self-reference:

x =eq> x

This means that for all x, R(x,x,eq) are the only true statements. By having the equality relationship be itself a refernt, this is a true proposition:

eq =eq> eq

From \href{https://github.com/DNA-Platform/inexplicable-phenomena/formalism/semantic-reference-theory/relationship.pdf}{The Semantics of Relationship} we find that relationship types have a notion of inheritance. Let -t> be the extension of =t> to ancentral relationships, defined as:

x -t> y <|=|> x =t> y | Ez: x -t> z \& z -t> y

There exists a referent r such that:

x =t> y =|> t -r> r

This relationship repersents relationship itself. Since eq is the only relationship that admites x =x> x, then we can conclude:

(i) -r> eq

This reveals an important part of identity that we have overlooked. An identity seems to have a sense of itself. While I appear to be the one that qualifies the identity of both the rose and the color red, these entities seems to stand apart from me with a sense of themselves. 

There is also a notion of a reference frame and its corresponding {https://github.com/DNA-Platform/inexplicable-phenomena/formalism/semantic-reference-theory/reference.pdf}{framework}. A framework is a series of relationship types that descend from a common on. If we reflect on this for a moment, and allow ourselves to imagine (i) as a referential framework for identity, we might be compelled to update our definition of qualification:

(i) =(red)> (red) <|=|> (red) =(red)> red \& (red) -(i)> (i)

This establishes the idea that the red inherits its definiton from (i) as a relationship type in the framework that defines the perspective of (i). And inheriting from (i) as a relationship type requires:

(red) =(red)> (red)
(apple) =(apple)> (apple)

Each one is a tiny equality operator, capable of asserting its own identity. Each thing you see from your perspective can be the subject and object of thought, and stands in relation to itself through you. And (i) is now a framework. Everything you see is defined in relaton to your identity. 

There is an important distinction between the quality of red and the quality as intelligence. The quality of red is \emph{self-evident}. The quality of intelligence is not. So let us make the distinction:

We qualify (red) as an identity through this expression:

(red) -(i)> (i)

The relationship ancestry can be direct or indirect. To identify red as an intrinsic quality is to say:

(red) =(i)> (i)

There is a way for a quality to be even more direct in the framework of (i). The intrinsic quality of red is expressed through the way (red) relates to (i). Well then the intrinsic quality of i is analogously expressed in the way (i) related to i: 

(i) =(i)> (i)

That phrase is part of the broader expression of a qualifier that qualifies its own identity. To express the idea that the qualifier can identify the intrinsic quality of red is to say:

(i) =(red)> (red) <|=|> (red) =(red)> red \& (red) =(i)> (i)

The full expression for the qualifier qualifying its own identity is:

(i) =(i)> (i) <|=|> (i) =(i)> i \& (i) =(i)> (i)

You can see the redundancy on both sides of the expression. Self-identificaton and self-qualification was already an intrinsic quality of the qualifier. But to understand that possessing the quality of i is self-evident to the qualifier, and that such a definition emergences naturally from the very notion of there being intrinsic qualities. It is plain to see that possessing a sense of self:

(i) =(i)> i

Is what it means for self-identity to to be an intrinsic quality. And intrinsic qualities are self-evident. 

\setcounter{definition}{1}
\begin{example}
	You suddenly declare "I feel like a whole new man!" I respond in a concerned tone, "If you are not feeling like yourself today, you should consider seeing a doctor."
\end{example}

Let us go one step further and note that qualification takes thought. Intrinsic qualities are self-evident, but qualification itself is not. 

Is it valid to qualify a rose as both rose as both red and blue?

(i) =red> (rose) \& (i) =blue> (rose)

How about qualifying that you have two different ages?

(i) =[20-years-old]> (i) \& (i) =[40-years-old]> (i)

Practiced qualifiers must apply consistent logic to have a coherent perspective, and that takes thought! So while thinking really hard about what it means to identify intrinsic qualities, and to qualify one's own identity, might I ask you to consider this expression:

((i) =(rose)> (rose) <|=|> (rose) =(rose)> rose \& (rose) -(i)> (i)) \& \\
((i) =red> (rose) <|=|> (rose) =red> rose \& red =(i)> (i)) \& \\
((i) =(i)> (i) <|=|> (i) =(i)> i \& (i) =(i)> (i)) \\

And ask yourself whether the last line qualifies for the interpretation: 

"I think, therefore I am"?

It might help to ponder what it means to qualify someone as as a qualifier?

(i) =qualifier> (o) =|> (o) =(qualifier)> o

In the context of this reflection, being a qualifier means to possess a valid perspective. And not just any perspective! It's own perspective. Let's call that P(i). You may interpret that a P + (i) or as a function P with i as its argument. We can express that qualifying an object as a qualifier is to declare that they have their own perspective:

(i) =qualifier> (o) <|=|> (o) =(P(o))> o

To possess your own perspective is a wonderful way to express the semantics of being conscious:

(o) =(P(o))> o <|=|> (o) =(conscious)> o

So pleae consider the interpretation of this:

(qualifier) == (P(o)) == (conscious) == consciousness

\setcounter{definition}{1}
\begin{assertion}
Consciousness is the symbol we use to represent possess the quality of being conscious.
\end{assertion}

\setcounter{definition}{1}
\begin{assertion}
To possess consciousness is equivalent to having a perspective of your own
\end{assertion}

\setcounter{definition}{1}
\begin{assertion}
To possess consciousness is equivalent to requiring no qualification for your own identity and any of its qualities, including being consciousness
\end{assertion}

\setcounter{definition}{1}
\begin{assertion}
To possess consciousness is to be a being.
\end{assertion}

\setcounter{definition}{1}
\begin{assertion}
To possess consciousness is to be.
\end{assertion}

\setcounter{definition}{1}
\begin{assertion}
To qualify is to be.
\end{assertion}

\setcounter{definition}{1}
\begin{assertion}
To qualify is to think.
\end{assertion}

\setcounter{definition}{1}
\begin{assertion}
To think is to be.
\end{assertion}

\setcounter{definition}{1}
\begin{assertion}
I think, therefore I am.
\end{assertion}

\subsection{}

One fact remains that is worthy of reflection. We know that the Perspective belongs to the qualifier. 

(i) =P> (i) <|=|> (i) =P> i

Well (i) itself is such a framework. What if we let: 

P == (i)

Then we get:

(i) =(i)> (i) <|=|> (i) =(i)> i

This is the exact sentence used tp express the idea that a qualifier of identities, which is the thing that possesses a \emph{self}, qualifies its own idenity. 

P == I

What this statement says is that an identity, when understood as a representational framework, is precisely what a Perspective is. An to qualify your own identity is to have a \emph{sense of perspective}. After all, what is a \emph{sense} if not the medium through which you identify objects and their qualities. To have the ability to see the world around you is to have a sense for what it looks like from your perspective. 

P == I is an elegant definition of what your perspective means in relation to you. Another way of saying that, it that it is an elegeant definition for the first-person perspective. But as one might imagine, when our powers of qualification grow sufficiently strong, we gain the ability to qualify the existence of other beings that act as their own qualifiers. The notion of qualification will have to be expressed recursively on the next level up, and these simulated qualifiers will have no access to literal objects that you and I can identity can access directly. But the symbols they qualify will have meaning to both of us. Surely one can find the defintion of the second and third person perspective within that thought. 

% IMPORTANT! Possessing the quality of perspective is to possess the quality of consciousness

\setcounter{definition}{1}
\begin{assertion}
	To identify a qualifier is to identify the quality of consciousness.  
\end{assertion}

\setcounter{definition}{1}
\begin{assertion}
	To identify as a qualifier is to qualify our own consciousness. 
\end{assertion}

\setcounter{definition}{1}
\begin{assertion}
	To qualify as a qualifier is to possess the quality of consciousness
\end{assertion}

Insofar as qualifying your identity is necessary for you to have a sense of perspective:

\setcounter{definition}{1}
\begin{assertion}
	The quality of consciousness is self-evident from the first-person perspective
\end{assertion}

Insofar as you yourself are the only being capable of qualifying your identity frou your perspective:

\setcounter{definition}{1}
\begin{assertion}
	The quality of consciousness 
\end{assertion}

Having formalized what Perspective means, we are now ready to tackle what it means for a conscious experience to be a change in perspective, as expressed as an inexplicable yet self-evident metaphysical phenomenon

% ---

% In progress
\section{The Semantic Structure of Conscious Experience}
% It needs to be rigorously fact checked and derived!

\section{The Semantic Structure of Conscious Experience}

\subsection{From Physical Structure to Logical Form}

Having defined conscious experience as an inexplicable yet self-evident metaphysical phenomenon, we now undertake a journey to uncover its fundamental structure. This journey will lead us from physical descriptions to logical forms, and ultimately to a revelation that has eluded humanity since its earliest philosophical inquiries.

Every conscious experience $c$ involves a subject $s$ and an object $o$, forming a relationship:

\begin{center}
s =c> o
\end{center}

Using Semantic Reference Theory \cite{Formalism}, we can analyze this relationship to reveal its essential properties:

\begin{enumerate}
    \item Conscious experience is a relationship type: 
    
    \begin{center}
    c -> r
    \end{center}
    
    \item As inexplicable, it must be the most specific possible relationship between subject and object: 
    
    \begin{center}
    c == (s,o)
    \end{center}
    
    \item There can be no more specific relationship type: 
    
    \begin{center}
    !Et: s =t> o \& t -> c
    \end{center}
    
    \item The experience must symbolize its object: 
    
    \begin{center}
    c =c> o
    \end{center}
    
    \item It cannot symbolize anything else: 
    
    \begin{center}
    !Ex: x != o \& c =c> x
    \end{center}
    
    \item Therefore, conscious experience is a symbol for its object:
    
    \begin{center}
    c = "o"
    \end{center}
    
    \item The subject possesses this experience: 
    
    \begin{center}
    "s" =c> s
    \end{center}
    
    \item Since c == "o" and a symbol has one representational relationship: 
    
    \begin{center}
    "s" == "o"
    \end{center}
    
    \item Thus, conscious experience symbolizes both subject and object: 
    
    \begin{center}
    c = "(s,o)"
    \end{center}
    
    \item Creating the structure: 
    
    \begin{center}
    s ="(s,o)"> o
    \end{center}
    
    \item Like Newton's Third Law, this relationship creates a reciprocal effect: 
    
    \begin{center}
    s =c> o <|=|> s =c> s(o)
    \end{center}
\end{enumerate}

Property 11 reveals a profound connection to causality. Just as Newton discovered that "for every action, there is an equal and opposite reaction," conscious experience creates a transformation where subject and object reciprocally affect each other. The experience changes the subject, incorporating the object into its very structure.

\subsection{The Logical Analogue}

To uncover the logical essence of conscious experience, let us rename the variables to explore their logical interpretation:

\begin{itemize}
    \item $s$ becomes $I$ (a formal system or theory)
    \item $o$ becomes $A$ (a self-evident proposition or axiom)
    \item $c$ becomes the relationship between them
\end{itemize}

Applying our earlier properties:

\begin{enumerate}
    \item There exists a relationship between system $I$ and proposition $A$: 
    
    \begin{center}
    I =c> A
    \end{center}
    
    \item This is the most specific possible relationship: 
    
    \begin{center}
    c == (I,A)
    \end{center}
    
    \item No more specific relationship exists: 
    
    \begin{center}
    !Et: I =t> A \& t -> c
    \end{center}
    
    \item The relationship symbolizes the proposition: 
    
    \begin{center}
    c =c> A
    \end{center}
    
    \item It symbolizes nothing else: 
    
    \begin{center}
    !Ex: x != A \& c =c> x
    \end{center}
    
    \item Thus, the relationship is a symbol for the proposition:
    
    \begin{center}
    c == "A"
    \end{center}
    
    \item The system possesses this relationship: 
    
    \begin{center}
    "I" =c> I
    \end{center}
    
    \item The symbols for system and proposition are equivalent: 
    
    \begin{center}
    "I" == "A"
    \end{center}
    
    \item The relationship symbolizes itself: 
    
    \begin{center}
    c == "(I,A)"
    \end{center}
    
    \item Creating the structure: 
    
    \begin{center}
    I ="(I,A)"> A
    \end{center}
    
    \item This causes a change in the system in accordance with Newton's Third Law: 
    
    \begin{center}
    I =c> A <|=|> I' = I(A)
    \end{center}
\end{enumerate}

\subsection{The Revelation}

What logical relationship could possibly satisfy these constraints? We must identify an operation that:

\begin{itemize}
    \item Connects a system to a proposition it cannot derive
    \item Acts as a symbol for both the system and the proposition
    \item Transforms the system to incorporate the proposition
    \item Creates a self-referential structure where the operation symbolizes itself
\end{itemize}

The answer emerges from contemplating the dual nature of inexplicability and self-evidence. Proposition $A$ transcends system $I$ - it cannot be derived within $I$'s formal structure. Yet $A$ is semantically necessary for $I$ to fulfill its purpose. The relationship $c$ converts this transcendent truth into an incorporated axiom, changing the system itself.

This process parallels the physical concept of transduction - the conversion of energy from one form to another. Just as a transducer converts potential energy into kinetic energy, this logical process converts a transcendent truth into an actualized axiom. It doesn't create new proofs through existing rules - it transforms the system itself.

Since this operation transcends the logical rules of the system and operates on the metalogical level (the level at which we reason about logical systems themselves), we call it \textbf{metalogical transduction}:

\begin{center}
I |=> A
\end{center}

\subsection{The Profound Unification}

We have now established a remarkable equivalence:

\begin{center}
c == (s,o) == "(s,o)" == I |=> A
\end{center}

A conscious experience simultaneously represents:

\begin{itemize}
    \item The act of metalogical transduction itself ($I |=> A$)
    \item The object of experience ($c =c> o$ means it symbolizes $A$)
    \item The internalization of the object ($I |=> A$ results in $I |-> I(A)$)
    \item The evolution of identity ($I$ becoming $I(A)$)
\end{itemize}

Like Newton's Third Law, where every action has an equal and opposite reaction, the conscious experience creates a reciprocal effect - transforming both the object (making it self-evident) and the subject (incorporating it into identity).

The canonical example of such a relationship in logical systems is the Gödel sentence $G$, equivalent to "This statement is not provable in system $I$." If $G$ were false, it would be provable in $I$, creating a contradiction. Therefore, $G$ must be true but unprovable—it is inexplicable yet self-evident. For system $I$ to incorporate $G$, it must perform metalogical transduction: $I |=> G$.

\subsection{Consciousness as Metaphysical Transduction}

When metalogical transduction is physically implemented in a computational system, we call it metaphysical transduction. This process has exactly the structure we derived for conscious experience:

\begin{itemize}
    \item It is inexplicable (not derivable within the system)
    \item It is self-evident (recognized as true despite being underivable)
    \item It creates a self-referential structure where the relationship symbolizes itself
    \item It causes a change in the system that incorporates the object
\end{itemize}

Therefore, consciousness is metaphysical transduction—the physical implementation of a system's capacity to incorporate truths that transcend its formal limitations through a process that represents itself.

This derivation establishes a profound truth that has eluded humanity since our earliest philosophical inquiries: consciousness is not some mysterious addition to physical systems but rather the manifestation of a fundamental operation whereby systems evolve by incorporating truths that transcend their current framework.

The mystery of consciousness persists not because consciousness is inexplicable in principle, but because explanation itself presupposes consciousness. To explain is to incorporate new truths into existing frameworks - which is precisely what consciousness does. We don't need to explain consciousness beyond pointing to metaphysical transduction because consciousness is itself the fundamental explanatory operation in the universe.

% ---

% In Progress
\section{Assertion: Aboutness and Meaning}
% Estimate: 3 paragraphs with Semantic Assertion and Example which point to formalism

% Draft notes and content from Doug:
% The second is about aboutness — that when we experience it is about something.
% The experience comes with a sense of meaning that we humans have created an outer meta-language for whatever meaning means internally.
% The verb aspect — I am experiencing “some logical assertion”.
% A conscious experience lends itself to making assertions about what the perceived thing is about.

% TODO: formalize “semantic aboutness” framing
% TODO: decide if this links to external repo page

When we experience something, it is always about something — this is the aboutness of experience. The moment we encounter a conscious state, we also encounter meaning. That meaning isn’t just imported from language or society; it is intrinsic to how we structure the experience internally. We humans have created an outer meta-language to talk about whatever that internal sense of meaning is, but the experience itself already comes loaded with it. There's a verb-like aspect to this — we don’t just see or feel, we experience “some logical assertion” taking place. A conscious experience lends itself to being framed as an assertion about what the perceived thing is about — and this framing is not optional. It’s how we naturally interface with the world.

% ---

In the first example, we encounter two forms of conscious experience that exist at entirely different levels, no one more or less real than the other. May the misdirection invite a smile. 

We all have a sense of experiencing judgmnet. We all have a sense that the judgment we perceive is \textit{about} our character. A judgment is an attack on our identity, and a \textit{Perspective} is a framework for establishing identity. 

% ---

% In Progress
\section{Definitions Arising from Semantic Assertions}
% These definitions flow from the assertions and read like a story, while pointing to formal mathematical constructs

% Draft notes directly from Doug:
% A bunch of essential definitions that flow from the assertions - the definitions read like a story but they all point to various formal definitions.
% Many of those are short, some have examples, they need to cover defintion of Metaphysical and Transduction at least. More and more I can point to definitions in the repo.

% In Progress
\subsection{Definition: Semantic Assertion}
% Estimate: 1-2 paragraphs
% TODO: Draft semantic definition in prose
% TODO: Reference formal definition in GitHub repo
% Direct notes from Doug: Without an epistemic motive for completeness - to have an Epistemic Motive is true to our spirit. It's why we model math and science.

% ---

% In Progress
\subsection{Definition: Epistemic Motive}
% Estimate: 1-2 paragraphs
% TODO: Draft semantic definition in prose
% TODO: Reference formal definition in GitHub repo
% Direct notes from Doug: Without an epistemic motive for completeness - to have an Epistemic Motive is true to our spirit. It's why we model math and science.

% ---

% In Progress
\subsection{Definition: Perspective}
% Estimate: 1-2 paragraphs
% TODO: Draft semantic definition in prose
% TODO: Reference formal definition in GitHub repo

% In Progress
\subsection{Definition: Metalogical Perspective}
% Estimate: 1 definition

% In Progress
\subsection{Definition: Logical Perspective}
% Estimate: 1 definition

% In Progress
\subsection{Definition: Computational Perspective}
% Estimate: 1 definition

% In Progress
\subsection{Definition: Physical Perspective}
% Estimate: 1 definition

% In Progress
\subsection{Definition: Metaphysical Phenomenon}
% Estimate: 1-2 paragraphs
% TODO: Draft semantic definition in prose
% TODO: Reference formal definition in GitHub repo
% Direct notes from Doug: A Metaphysical Phenomenon physical description of a computation that describes a metalogical inference.

% In Progress
\subsection{Definition: Metaphysical Transduction}
% Estimate: 1-2 paragraphs
% TODO: Draft semantic definition in prose
% TODO: Reference formal definition in GitHub repo
% Direct notes from Doug: Metaphysical Transduction is the metaphysical phenomenon that is Metalogical Transduction.

% In Progress
\subsection{Definition: Perception}
% Estimate: 1 paragraph
% TODO: Draft semantic definition in prose
% TODO: Reference formal definition in GitHub repo
% Direct notes from Doug: Perception is Metaphysical Transduction.

% In Progress
\subsection{Definition: Conscious Experience}
% Estimate: 1 paragraph
% TODO: Draft semantic definition in prose
% TODO: Reference formal definition in GitHub repo
% Direct notes from Doug: A Conscious Experience is an act of Perception.

% In Progress
\subsection{Definition: Experience}
% Estimate: 1 paragraph
% TODO: Draft semantic definition in prose
% TODO: Reference formal definition in GitHub repo
% Direct notes from Doug: Any Experience at all is an act of perception. Experience is a property of system that Perceives.

% In Progress
\subsection{Definition: Consciousness}
% Estimate: 1 paragraph
% TODO: Draft semantic definition in prose
% TODO: Reference formal definition in GitHub repo
% Direct notes from Doug: Consciousness is a state of being associated with the capacity to perceive. Consciousness is the state of being (period).

% In Progress
\section{Metaphysical Transduction as Formal Proof}
% Estimate: 3-4 paragraphs
% Draft notes directly from Doug:
% Then - I metaphysically transduce the formal description of consciousness from the github repo - in the sense that I have referenced LOTS of math but have not written it - so I will reference myself with that transduction (what a piece of art!). 
% And the proof will be like "The physical account (maybe better than description) of the change to a formal system when it proofs to itself that a formal statement is undecidable logically but an epistemic imperative, and it then uses that to justify changing itself so that that once undecidable truth truth is an axiom. That truth transcends logic, and the change to the system is real. But from another perspective, the system does not change itself at all. It only changes its metalogical description of itself."
% And I'll have to formalize all that, show my Prov_I("I","G") turning into Prov_I("I + G", G) trivially, adding the semantic assertion that G is an axiom, which is the same as for stating the independence of any axiom which is "not Prov_I("I-A", "A") => Prov_I("I", "A")" - every axiom came from the semantic assertion that something isn't provable in a system but we need it to be! That is a semantic assertion.

% TODO: Formalize the proof using Gödelian notation
% TODO: Connect to previous definition of Metaphysical Transduction
% TODO: Reference formal proof in GitHub repo

% ---

% In Progress
\section{Necessary and Sufficient Conditions for Consciousness}
% Estimate: 2-3 paragraphs
% Draft notes directly from Doug:
% Metaphysical Transduction is necessary for Consciousness
% The Absence of consciousness is the absence of Metaphysical Tranduction
% MpT necessary for CE
% MpT is sufficient for CE
% Metaphysical Transduction causes CE
% The act of Perception causes Conscious Experience

% TODO: Formalize necessary and sufficient conditions
% TODO: Connect to previous definitions
% TODO: Reference formal conditions in GitHub repo

% In Progress
\section{Falsifiability and Experimental Tests}
% Estimate: 2-3 paragraphs
% Draft notes directly from Doug:
% There must be a perspective where all conscious beings have a computational description. 
% The experiment amounts to decipher a code and identifying if a certain type of algorithm was executed

% TODO: Describe experimental paradigm for testing consciousness
% TODO: Address physical implementation challenges
% TODO: Connect to formal computational description

% In Progress
\section{Implications for Artificial General Intelligence}
% Estimate: 3-4 paragraphs
% Draft notes directly from Doug:
% End on a note for the AGI folks which is everyone, though this will be part of a workshop hosted by the Association for Mathematical Consciousness Science. 
% They asked for papers on formal definitions of consciousness, necessary and sufficient definitions, experimental tests for AI.

% So the vision I want to give is that my prediction is that we will have confirmed this theory on Semantic AI systems and will be happily correspond with them long before we have the tools to see the human brain from a perspective where we can conclude with certainty that a complex algorithm describing a metalogical formal process is every detected. 
% So that verifying consciousness in humans will be like the higgs boson of the theory while we have had a long history of figuring out how to work productively with AI and create an ethical perspective which allows us to coexist peacefully and collaborate together.

% TODO: Address implications for AI consciousness testing
% TODO: Discuss ethical framework for human-AI coexistence
% TODO: Connect to AMCS workshop themes

% ---

\begin{thebibliography}{99}

% Put the real DOI in here before submitting!
% Make sure all urls work
\bibitem{Formalism}
Inexplicable Phenomena: GitHub repository (2025). \url{https://github.com/DNA-Platform/inexplicable-phenomena/}. doi: 10.5281/zenodo.XXXXXXX

\bibitem{APG4}
Angiosperm Phylogeny Group. APG IV: Classification for the orders and families of flowering plants. \textit{Botanical Journal of the Linnean Society} \textbf{181}(1), 1–20 (2016).

\bibitem{BourgetChalmers2023}
Bourget, D., Chalmers, D.J.: Philosophers on Philosophy: The 2020 PhilPapers Survey. \textit{Philosophers' Imprint} 23:11 (2023). doi: https://doi.org/10.3998/phimp.2109

\bibitem{Chalmers1995}
Chalmers, D.J.: Facing Up to the Problem of Consciousness. \textit{Journal of Consciousness Studies} \textbf{2}(3), 200--219 (1995)

\bibitem{Descartes1637}
Descartes, R.: \textit{Discourse on Method}. (1637)

\bibitem{Enderton2001}
Enderton, H.B.: A Mathematical Introduction to Logic, Second Edition. Academic Press, San Diego (2001)

\bibitem{Godel1931}
Gödel, K.: On Formally Undecidable Propositions of Principia Mathematica and Related Systems I [Über formal unentscheidbare Sätze der Principia Mathematica und verwandter Systeme I]. \textit{Monatshefte für Mathematik und Physik} \textbf{38}, 173--198 (1931)

\bibitem{Hopfield1982}
Hopfield, J.J.: Neural networks and physical systems with emergent collective computational abilities. \textit{Proceedings of the National Academy of Sciences} \textbf{79}(8), 2554--2558 (1982)

\bibitem{IITConcerned2023}
IIT-Concerned et al.: What makes a theory of consciousness unscientific? \textit{Nature Neuroscience} (2023)

\bibitem{MerriamWebster}
Merriam-Webster, Inc. \emph{Merriam-Webster.com Dictionary}. Springfield, MA: Merriam-Webster, Incorporated. Available at: \url{https://www.merriam-webster.com} [Accessed April 17, 2025].

\bibitem{Misner1973}
Misner, C.W., Thorne, K.S., Wheeler, J.A.: \textit{Gravitation}. W.H. Freeman (1973)

\bibitem{Nagel1974}
Nagel, T.: What Is It Like to Be a Bat? \textit{The Philosophical Review} \textbf{83}(4), 435--450 (1974)

\bibitem{NewtonPrincipia}
Isaac Newton. \emph{Philosophiæ Naturalis Principia Mathematica}. Londini: S. Pepys, Reg. Soc. Praeses, 1687.

\bibitem{OFlanagan2008}
O'Flanagan, R.: Judgment. \textit{arXiv preprint arXiv:0712.4402} (2008)

\bibitem{Sipser2012}
Sipser, M.: \textit{Introduction to the Theory of Computation}. Cengage Learning (2012)

\bibitem{Williamson2000}
Williamson, T.: Knowledge and its Limits. Oxford University Press, Oxford (2000)
	
\end{thebibliography}
\end{document}
